\documentclass{article}
\usepackage{amsmath, amssymb}
\setlength{\parindent}{0pt}

\begin{document}

\textbf{CS 2050 --- Notes 08/31}

\bigskip

\textbf{Compound proposition.} Any expression formed from propositional variables using logical operators.

\medskip

\textbf{Example.} $p \land q$; \quad $(p \lor \lnot q) \to (\lnot p \land \lnot q)$

\bigskip

\textbf{Tautology.} A compound proposition that is always true.

\medskip

\textbf{Example.} $p \to p$ is a tautology.

\bigskip

\textbf{Contradiction.} A compound proposition that is always false.

\medskip

\textbf{Example.} $\lnot p \land p$ is a contradiction.

\bigskip

\textbf{Contingency.} A compound proposition that is neither a tautology nor a contradiction.

\bigskip

\textbf{De Morgan's Laws}
\begin{align*}
\lnot (p \land q) &\equiv \lnot p \lor \lnot q & &\text{De Morgan's 1st Law} \\
\lnot (p \lor q) &\equiv \lnot p \land \lnot q & &\text{De Morgan's 2nd Law}
\end{align*}

\bigskip

\textbf{Conditional-disjunction equivalence:} $p \to q \equiv \lnot p \lor q$

\medskip

\textbf{Contrapositive law:} $p \to q \equiv \lnot q \to \lnot p$

\medskip

\textbf{Definition of biconditional:} $p \leftrightarrow q \equiv (p \to q) \land (q \to p)$

\bigskip

\textbf{Double negation law:} $\lnot (\lnot p) \equiv p$

\medskip

\textbf{Identity Laws}
\begin{align*}
p \land T &\equiv p \\
p \lor F &\equiv p
\end{align*}

\medskip

\textbf{Domination Laws}
\begin{align*}
p \land F &\equiv F \\
p \lor T &\equiv T
\end{align*}

\medskip

\textbf{Idempotent Laws}
\begin{align*}
p \lor p &\equiv p \\
p \land p &\equiv p
\end{align*}

\medskip

\textbf{Commutative Laws}
\begin{align*}
p \lor q &\equiv q \lor p \\
p \land q &\equiv q \land p
\end{align*}

\medskip

\textbf{Associative Laws}
\begin{align*}
(p \lor q) \lor r &\equiv p \lor (q \lor r) \\
(p \land q) \land r &\equiv p \land (q \land r)
\end{align*}

\medskip

\textbf{Distributive Laws}
\begin{align*}
p \lor (q \land r) &\equiv (p \lor q) \land (p \lor r) \\
p \land (q \lor r) &\equiv (p \land q) \lor (p \land r)
\end{align*}

\medskip

\textbf{Absorption Laws}
\begin{align*}
p \lor (p \land q) &\equiv p \\
p \land (p \lor q) &\equiv p
\end{align*}

\medskip

\textbf{Negation Laws}
\begin{align*}
p \lor \lnot p &\equiv T \\
p \land \lnot p &\equiv F
\end{align*}

\bigskip

\textbf{Satisfiability.} A compound proposition is satisfiable if there is an assignment of truth values to its variables that makes it true.

\medskip

\begin{itemize}
\item Way to show: truth table --- a compound proposition is satisfiable if one T is present in the table.
\item Logically reducing the statement is possible as well.
\end{itemize}

\bigskip

\textbf{Predicate and logic}

\medskip

\textbf{Propositional logic}
\begin{itemize}
\item Druids are magic users.
\item $2 \ast 9 = 42$.
\end{itemize}

\medskip

\textbf{Propositional logic failed:} $x + y = 27$

\bigskip

\textbf{Example.} $x < 3$: $x$ is the subject, $3$ is the predicate. The statement ``$x < 3$'' can be represented by $P(x)$.
\[
P(0) = 0 < 3 = \text{true}, \qquad P(4) = 4 < 3 = \text{false}
\]

\medskip

\textbf{Example.} Let $D(x)$ denote ``someone in this class likes games.''
\[
D(\text{Jade}) \equiv \text{Jade likes games}
\]

\bigskip

\textbf{Multivariable}

\medskip

\textbf{Example.} Let $Q(x, y)$ denote $x = 3 - y$. What are the truth values of $Q(1, 2)$ and $Q(2, 3)$?
\begin{align*}
Q(1, 2) &= (1 = 3 - 2) = \text{true} \\
Q(2, 3) &= (2 = 3 - 3) = \text{false}
\end{align*}

\bigskip

$\forall$ is the universal quantifier.

\medskip

$\exists$ is the existential quantifier.

\medskip

Universe of discourse = domain.

\medskip

$\exists !$ is the uniqueness quantifier.

\bigskip

\textbf{Quantifiers over finite domain.} If the domain of a quantifier is finite, then a quantified statement can be converted to propositional logic:
\[
\forall x P(x) \equiv P(x_1) \land P(x_2) \land \cdots \land P(x_n)
\]
\[
\exists x P(x) \equiv P(x_1) \lor P(x_2) \lor \cdots \lor P(x_n)
\]

\end{document}
